\section{Data Requirements}

The model has two sets of parameters to be estimated: preference parameters $(\eta, \gamma, \rho_1, \rho_2)$ and human capital parameters $(\bar{w}_j, \alpha_j, \delta)$.
These would be identified via Simulated Method of Moments (SMM), matching the model's simulated moments to the following empirical targets:

\begin{itemize}
    \item Mean hours worked by gender and age group.
    \item The gender wage gap profile over the life cycle (ages 25--55).
    \item The within-household hours ratio at different stages of the career.
    \item The earnings elasticity of female labor supply with respect to the net-of-tax rate.
\end{itemize}

The natural data source is Danish administrative register data (Statistics Denmark's IDA database), which provides annual earnings, hours worked, household composition, and tax records for the full population.
Denmark is ideal for two reasons: the data quality is high, and Denmark abolished joint taxation in 1983, providing a quasi-experimental source of variation.
Estimation could exploit a difference-in-differences design comparing households that were differentially affected by the 1983 reform --- those where the secondary earner's income crossed a key bracket threshold under joint filing --- to pin down the behavioral parameters.

\newpage
\bibliographystyle{plainnat}
\bibliography{references}
